%! AP Statistics — Unit 1, Lesson 1.6 — Homework ANSWER KEY
\documentclass[10pt]{article}
\usepackage{apstats-article}
\usepackage{apstats-key}

%=============================================================================
\begin{document}

\pageheader{Unit 1, Lesson 1.6}{Homework --- Answer Key}
\begin{tabularx}{\linewidth}{X X X}
Name:\;\ans{Key} & Date:\; & Period:\;
\end{tabularx}

\vspace{0.1in}

% ════════════════════════════════════════════════════════════════════════════
% PART A
% ════════════════════════════════════════════════════════════════════════════
\begin{blurbbox}[Part A \quad Identify the Shape --- Key]{navylight}
\small

\begin{enumerate}[leftmargin=1.5em, itemsep=10pt]

  \item Shape: \ans{Skewed right.}
        Evidence: \ans{Most homes sell in the \$180k--\$350k range (left bulk),
        but luxury estates above \$2M pull a long tail to the right.}

  \item Shape: \ans{Skewed left.}
        Evidence: \ans{Most students scored 9--10 (right bulk), with only a
        few scoring below 7 --- a tail toward lower (left) values.}

  \item Shape: \ans{Symmetric (unimodal).}
        Evidence: \ans{The peak near 7 hours has roughly equal drop-off on
        both sides, forming a bell shape.}

  \item Shape: \ans{Uniform.}
        Evidence: \ans{All five response options were chosen by nearly equal
        numbers of respondents, so all bars are roughly the same height.}

\end{enumerate}

\begin{teachernote}
Part A is a shape identification exercise. Award credit for the correct shape
\textit{and} a reasonable piece of evidence from the description. Reject answers
that name the shape without any justification.
\end{teachernote}

\end{blurbbox}

\vspace{0.14in}

% ════════════════════════════════════════════════════════════════════════════
% PART B
% ════════════════════════════════════════════════════════════════════════════
\begin{blurbbox}[Part B \quad SOCS Description --- Key]{greenacc}
\small

\begin{enumerate}[leftmargin=1.5em, itemsep=10pt]

  \item \ans{The distribution of daily physical activity time for these
        45 middle school students is skewed right, with most students being
        active between 20 and 60 minutes per day and a long tail extending
        toward higher activity times. [S] \quad One student reported 150 minutes
        of physical activity per day, a potential outlier that falls far above
        the rest of the group. [O] \quad The center of the distribution is
        approximately 38 minutes per day, representing a typical level of daily
        physical activity for a middle school student in this sample. [C] \quad
        Daily physical activity times range from 5 to 150 minutes, a spread of
        145 minutes. [S]}

  \begin{teachernote}
  Rubric: 1 pt each for S (skewed right, evidence), O (outlier identified with
  approximate value), C (center with units and context), S (range with units).
  Deduct for missing context in any sentence.
  \end{teachernote}

  \item \ans{Removing the 150-minute student would have the following effects:\\
        \textit{Shape:} Would remain skewed right, but less so --- the tail
        would shorten now that the extreme value is gone.\\
        \textit{Center:} Would decrease slightly (the mean would drop more than
        the median; the median would be approximately unchanged).\\
        \textit{Spread:} Would decrease noticeably --- the maximum would drop
        from 150 to around 90--120 minutes, reducing the range.}

  \item \ans{For a skewed-right distribution, the few unusually large values
        (the long right tail) pull the mean toward higher values more than they
        pull the median, so the mean is typically \textit{greater than} the
        median. Here, the mean daily activity time is likely greater than 38
        minutes (the approximate median), because the student at 150 minutes
        inflates the mean.}

\end{enumerate}

\end{blurbbox}

\vspace{0.14in}

% ════════════════════════════════════════════════════════════════════════════
% PART C
% ════════════════════════════════════════════════════════════════════════════
\begin{blurbbox}[Part C \quad AP-Style Free Response --- Key]{redacc}
\small

\begin{enumerate}[leftmargin=1.5em, itemsep=10pt]

  \item \ans{The distribution of fuel efficiency for these 60 cars is roughly
        symmetric and unimodal, with no obvious skew. [S] \quad There appears
        to be one potential outlier at 47 mpg, which falls above the main body
        of the distribution; no cars appear to be unusual at the low end. [O]
        \quad The center of the distribution is approximately 32 miles per
        gallon, representing the typical fuel efficiency for a car in this
        sample. [C] \quad Fuel efficiency values range from 18 to 47 miles per
        gallon, a spread of 29 miles per gallon. [S]}

  \item \ans{Adding 5 cars at 14 mpg would change SOCS as follows:\\
        \textit{Shape:} Would shift to skewed left --- the new cars create a
        tail toward lower mpg values.\\
        \textit{Outliers:} The five 14-mpg cars would become potential outliers,
        falling well below the existing minimum of 18 mpg.\\
        \textit{Center:} Would decrease (both mean and median would shift lower,
        with the mean decreasing more).\\
        \textit{Spread:} Would increase --- the minimum drops from 18 to 14 mpg,
        widening the range.}

  \item \ans{Yes, the classmate is correct for roughly symmetric distributions.
        When a distribution is symmetric, the mean and median are approximately
        equal because the symmetric values on either side of center cancel out
        and do not pull the mean in either direction. The 47-mpg outlier slightly
        inflates the mean, so the mean may be just above the median, but the
        difference would be small.}

\end{enumerate}

\end{blurbbox}

\end{document}
